\begin{ledgroupsized}[r]{120mm}
\footnotesize
\pstart
\noindent\textbf{Überlieferung:}
\pend
\end{ledgroupsized}
\begin{ledgroupsized}[r]{114mm}
\footnotesize
\pstart \parindent -6mm
\makebox[6mm][l]{\textit{L}}%
Aufzeichnung:
LH~XXXVII~5~Bl.~63.
Ein Zettel (ca.~18~x~5,5~cm.).
Papiererhaltungsmaßnahmen;
oberer und rechter Rand beschnitten.
Eine Seite auf Bl.~63~r\textsuperscript{o}; Bl.~63~v\textsuperscript{o} leer.
\pend
\end{ledgroupsized}
\selectlanguage{latin}
\vspace{8mm}
\pstart%
\normalsize%
\noindent%
\href{\Scan{LH_37_05_063r}}%
{\lbrack63~r\textsuperscript{o}\rbrack}\
(1\textsuperscript{mo}) $av+bx=ay
\edtext{+bz$. ((1\textsuperscript{mo})) seu $v-y,:,-x+z\lbrack=\rbrack b:a$. (2\textsuperscript{do}) $v-x=-y+z$. Ex 2}{%
\lemma{}\Bfootnote{%
$+bz$. 
\bars{((1\textsuperscript{mo})) seu 
\bars{$v-y,:,-x+z : b:a$  \textit{ändert~Hrsg.}}
\bars{seu differenti \textit{streicht Hrsg.}}
\textit{erg.}}
.~(2\textsuperscript{do})
\textit{(1)}~$v-y$ %
\textit{(2)}~$v-x=-y+z$.
\textit{(a)}~$z=$
\textit{(b)}~3\textsuperscript{tio} $z=v-x+y$. 4\textsuperscript{to} 
\textit{(c)}~Ex 2~\textit{L}}} 
fit (3\textsuperscript{tio}) $z=v-x+y$.
Ex 3 et 1 fit (4\textsuperscript{to}) $av+bx=ay+bv-bx+by$. 
Et hinc fit 
\edtext{(5\textsuperscript{to}) $(a-b)v}{\lemma{(5\textsuperscript{to})}\Bfootnote{%
\textit{(1)}~$av+bx$ %
\textit{(2)}~$\overline{a-b}\; v+2bx=\overline{a+b}\; y$ %
\textit{(3)}~$(a-b)v\protect\vphantom()+$ %
\textit{(4)}~$(a-b)v$~\textit{L}}} 
+2bx=(a+b)y$ et pari jure 
\edtext{fit (6\textsuperscript{to})}{\lemma{fit}\Bfootnote{%
\textit{(1)}~$2by$ %
\textit{(2)}~(6\textsuperscript{to})~\textit{L}}} 
$(a-b)z+2by=\edlabel{37_05_063_1a}(a+b)x$.%
\edtext{}{{\xxref{37_05_063_1a}{37_05_063_1b}}\lemma{$(a+b)x$.}\Bfootnote{%
\textit{(1)}~Fit %
\textit{(2)}~Aestimemus~\textit{L}}}
\pend
\pstart
Aestimemus\edlabel{37_05_063_1b} 
vires impeditivas\indexS{vis impeditiva}
\edtext{seu progressiones.\indexS{progressio}}{\lemma{}\Bfootnote{seu progressiones \textit{erg.~L}}} 
\edtext{(7\textsuperscript{mo})}{\lemma{}\Bfootnote{(7\textsuperscript{mo}) \textit{erg.~L}}} 
$av=11$\lbrack,\rbrack\
\edtext{(8\textsuperscript{vo})}{%
\lemma{}%
\Bfootnote{%
(8\textsuperscript{vo}) %
\textit{erg.~L}}}
$bx=21$\lbrack,\rbrack\
(9\textsuperscript{no}) 
\edtext{$ay=12$\lbrack,\rbrack\ %
\edtext{(10\textsuperscript{mo})}{\lemma{(10\textsuperscript{mo})}%
\Cfootnote{Diese Zählung wird von Leibniz zweimal verwendet.}}
$bz=22$. Ergo (10\textsuperscript{mo})}{%
\lemma{$ay=12$}%
\Bfootnote{%
\bars{(10\textsuperscript{mo}) \textit{erg.}} $bz=22$. Ergo
\textit{(1)}~$v=11:a$. \textit{x} %
\textit{(2)}~(10\textsuperscript{mo})~\textit{L}}}
$v=11:a$\lbrack,\rbrack\
(11\textsuperscript{mo}) $x=21:b$\lbrack,\rbrack\
(12\textsuperscript{o}) $y=12:a$\lbrack,\rbrack\
(13\textsuperscript{o}) $z=22:b$.
Quos 
\edtext{valores ex articulis 10, 11, 12, 13,}{\lemma{valores}\Bfootnote{%
\textit{(1)}~in aeqq. %
\textit{(2)}~ex %
\textit{(a)}~aeqq. 7, 8, 9, 10 %
\textit{(b)}~articulis %
\textit{(aa)}~7, 8, 9, 10 %
\textit{(bb)}~10, 11, 12, 13,~\textit{L}}} 
substituendo 
\edtext{in 1, et 2.\ utique}{\lemma{in}\Bfootnote{%
\bars{aeqq. \textit{gestr.}} 1, et 2.\
\textit{(1)}~fiet ex 1.\ %
\textit{(2)}~\bars{utique ex 1.\ \textit{streicht~Hrsg.}} %
\textit{(3)}~fiet ex 1.\ %
\textit{(4)}~utique~\textit{L}}} 
ex 1.\ fiet
\edtext{(14\textsuperscript{o})}{\lemma{}\Bfootnote{(14\textsuperscript{o}) \textit{erg.~L}}} 
$11+21=12+22$ et 
\edtext{ex 2 fiet (15\textsuperscript{to})}{\lemma{ex 2 fiet}\Bfootnote{%
\textit{(1)}~15\textsuperscript{to} %
\textit{(2)}~(15\textsuperscript{o}) %
\textit{(3)}~(15\textsuperscript{to})~\textit{L}}} 
$(11:a)-(21:b)=
\edtext{-(12:a)+\lbrack(22:b)\rbrack$ vel (16\textsuperscript{to})}{\lemma{}\Bfootnote{$-(12:a)$ 
\bars{$+22:b$ \textit{ändert~Hrsg.}}
\textit{(1)}~adeoque fiet
\textit{(2)}~vel
\textit{(a)}~$(11+12)$
\textit{(b)}~(16\textsuperscript{to})~\textit{L}}} 
$11+12,:a=21+22,:b$ vel (17\textsuperscript{mo}) 
$11+12,:21+22=
\edtext{a:b$. Ergo (18\textsuperscript{vo}) $11b-21a=-12b+22a$. Ejusdem}{\lemma{$a:b$}\Bfootnote{%
\textit{(1)}~seu differentiae vel su %
\textit{(2)}~. Si ad easdem partes t %
\textit{(3)}~. Si ad %
\textit{(4)}~\bars{Ergo (18\textsuperscript{vo}) $11b-21a=-12b+22a$. \textit{erg.}} %
Ejusdem~\textit{L}}} 
corporis si ad easdem partes
\edtext{tendat summae,\indexS{summa virium impeditivarum ejusdem corporis}}{\lemma{tendat}\Bfootnote{%
\textit{(1)}~differentiae %
\textit{(2)}~summae,~\textit{L}}}
si ad 
\edtext{contrarias differentiae\indexS{differentia virium impeditivarum ejusdem corporis}}{\lemma{contrarias}\Bfootnote{%
\textit{(1)}~summae %
\textit{(2)}~differentia %
\textit{(3)}~differentiae~\textit{L}}} 
virium impeditivarum\indexS{vis impeditiva} sunt ut corpora; 
\edtext{vel generalius summae}{\lemma{vel}\Bfootnote{\bars{generalius \textit{erg.}} %
\textit{(1)}~differentiae %
\textit{(2)}~summae~\textit{L}}} 
progressionum\indexS{progressio} 
\edtext{ejusdem corporis\indexS{summa progressionum ejusdem corporis}}{\lemma{}\Bfootnote{ejusdem corporis \textit{erg.~L}}} 
ante ictum et post ictum\indexS{ictus} 
\ledpb%
sunt ut 
\edtext{corpora. Ex 1 summa progressionis}{\lemma{corpora}\Bfootnote{%
\textit{(1)}~, quia differentia virium im %
\textit{(2)}~, quia
\textit{(3)}~, generalius loquendo, %
\textit{(4)}~, quod ita generale est. %
\textit{(5)}~. Jam tollatur valor ipsius 22 %
\textit{(a)}~ex 13 %
\textit{(b)}~per 14, fiet (18\textsuperscript{vo}) $22=11+21-12$ ergo ex 17 fiet (19\textsuperscript{no}) $11-12,:, 21-11 - 21+12$ %
\textit{(6)}~. Memorabile est \bars{(18\textsuperscript{vo}) \textit{erg.}} $11-12=-21+22$ et summa progress %
\textit{(7)}~Aeq. 1 %
\textit{(8)}~. Ex 1 summa progressionis~\textit{L}}} 
ante ictum\indexS{summa progressionis ante ictum} aequalis summae progressionis post ictum.\indexS{summa progressionis post ictum} 
Ex 2 differentia velocitatum progrediendi\indexS{differentia velocitatum progrediendi} ante ictum\indexS{ictus}
aequalis 
differentiae velocitatum progrediendi\indexS{differentia velocitatum progrediendi} post ictum.\indexS{ictus} 
Seu velocitas respectiva\indexS{velocitas respectiva} eadem. 
Posses quantitatem progressus\indexS{quantitas progressus} et vim impeditivam\indexS{vis impeditiva}
pro eodem sumere.
\pend 
